% Chapter — Spherical-harmonic gravity (FR-5). Follows the mandatory FR-29
% six-section template: purpose; assumptions and domain of validity;
% derivation; discretization and implementation; validation evidence;
% references. The equation labels eq:gravity:potential, eq:gravity:pines,
% eq:gravity:diag, eq:gravity:subdiag, eq:gravity:column, eq:gravity:deriv,
% eq:gravity:rmim, eq:gravity:sums, and eq:gravity:accel are echoed verbatim
% in comments in cpp/src/models/gravity.cpp and
% cpp/include/star/models/gravity.hpp (FR-29 equation-label-to-code
% traceability); renaming them breaks that trace.
\chapter{Spherical-Harmonic Gravity}
\label{ch:gravity}

\section{Purpose}
\label{sec:gravity:purpose}

This chapter documents the harmonic gravity model (FR-5): evaluation of
the gravitational acceleration of an extended central body from a fully
normalized spherical-harmonic coefficient set --- Earth EGM2008, Moon
GRGM1200A, Mars MRO120F --- in the body-fixed frame, using the Pines
formulation, which is regular at the poles. The model supports a runtime
truncation degree and order and the three FR-5 fidelity tiers
(point-mass, $J_2$-only, full $n \times m$). Field data arrives either
from in-memory arrays (tests) or from the SRGRAV~v1 binary container
written by \texttt{star data fetch}
(\texttt{docs/formats/srgrav\_v1.md}); the returned acceleration always
includes the central (degree-0) term. The module is implemented in
\texttt{cpp/src/models/gravity.cpp} with its interface in
\texttt{cpp/include/star/models/gravity.hpp}; body-frame orientation is
the frames chapter's responsibility (Chapter~\ref{ch:frames}).

\section{Assumptions and domain of validity}
\label{sec:gravity:domain}

Assumptions:
\begin{enumerate}
  \item \textbf{Static, rigid field.} The coefficient set is constant:
    solid and ocean tides, atmospheric loading, and secular
    $\bar{C}_{2,0}$ drift are unmodeled. Their effect on the low-degree
    coefficients is at the $10^{-8}$--$10^{-9}$ level of
    $\bar{C}_{2,0}$~\cite{pavlis2012}, far below the mission-analysis
    fidelity this simulator targets.
  \item \textbf{Fully normalized inputs.} Coefficients are $4\pi$
    (geodesy) fully normalized $\bar{C}_{nm}, \bar{S}_{nm}$; the data
    pipeline rejects any source that declares another normalization.
    Each field's $\mu = GM$ and reference radius $R$ come from the same
    source header as its coefficients --- the triple is self-consistent
    and is never mixed with constants from elsewhere. In particular
    EGM2008's $\mu = \qty{3.986004415e14}{\metre\cubed\per\second\squared}$
    differs deliberately from the IERS value in
    \texttt{constants.hpp}~\cite{iers2010} used by the two-body
    placeholder model.
  \item \textbf{Permanent-tide convention documented, not applied.} The
    tide system of each source (EGM2008 tide-free per its header;
    MRO120F tide-free per its PDS label; GRGM1200A not stated by its
    label, recorded as unknown) is carried through the SRGRAV header for
    cross-tool bookkeeping. It never enters evaluation; it only affects
    which $\bar{C}_{2,0}$ convention a comparison tool must be fed.
  \item \textbf{Exterior evaluation.} The exterior series converges for
    $r$ greater than the Brillouin sphere radius (here the reference
    radius $R$); all validation states satisfy $r > R$. Below $R$ the
    series may diverge and results are not trustworthy.
  \item \textbf{Body-fixed input, body-fixed output.} Position and
    acceleration are resolved in the source model's body-fixed frame
    (ITRS-aligned for EGM2008, lunar principal axes for GRGM1200A, IAU
    Mars body-fixed for MRO120F). Rotation into the integration frame is
    composed from Chapter~\ref{ch:frames} and is outside this module.
\end{enumerate}

Domain of validity: degrees up to the shipped maxima ($70\times70$
Earth, $120\times120$ Moon, $80\times80$ Mars; FR-5). The normalized
recursions of Section~\ref{sec:gravity:derivation} are the
forward-stable column recursions recommended for nonsingular
geopotential models~\cite{lundberg1988}; the normalized Helmholtz values
$\bar{A}_{nm}(u)$ grow to at most $\sim 10^{23}$ near $|u| = 1$ at
degree 120 while their multipliers $r_m, i_m$ decay as
$\cos^m\varphi$, so every summed term is bounded by the corresponding
normalized-Legendre magnitude $O(\sqrt{2n+1})$ and binary64 overflow is
excluded by orders of magnitude at the shipped degrees.

Out-of-domain behavior, matching the code exactly: a requested
truncation degree or order beyond the stored band throws
\texttt{std::invalid\_argument} (the file carries no information above
its band; fidelity is never silently degraded); evaluation at $r = 0$
throws \texttt{std::domain\_error}; the $J_2$-only tier on a field of
degree $< 2$ throws \texttt{std::invalid\_argument}; a malformed SRGRAV
file (magic, major version, size, tide code) is rejected by the loader
with \texttt{std::runtime\_error} naming the defect. Evaluation at
$R_{\text{surface}} < r < R$ (possible for the Moon and Mars, whose
reference radii exceed parts of the physical surface) proceeds without
warning; assumption 4 bounds its trustworthiness.

\section{Derivation}
\label{sec:gravity:derivation}

\subsection{Normalized harmonic potential}

The exterior gravitational potential of an extended body, in body-fixed
spherical coordinates (radius $r$, geocentric latitude $\varphi$,
longitude $\lambda$), is the standard harmonic
expansion~\cite{vallado2013, pavlis2012}
\begin{equation}
  V = \frac{\mu}{r} \sum_{n=0}^{N} \left(\frac{R}{r}\right)^{\!n}
      \sum_{m=0}^{\min(n,M)} \bar{P}_{nm}(\sin\varphi)
      \left[ \bar{C}_{nm} \cos m\lambda + \bar{S}_{nm} \sin m\lambda
      \right],
  \label{eq:gravity:potential}
\end{equation}
with $\bar{P}_{nm} = N_{nm} P_{nm}$ the $4\pi$ fully normalized
associated Legendre functions,
\begin{equation}
  N_{nm} = \sqrt{\frac{(2 - \delta_{0m})(2n+1)(n-m)!}{(n+m)!}},
  \label{eq:gravity:norm}
\end{equation}
and $\bar{C}_{nm} = C_{nm}/N_{nm}$, $\bar{S}_{nm} = S_{nm}/N_{nm}$ the
correspondingly normalized Stokes coefficients. The $n = 0$ term is
$\mu \bar{C}_{0,0} / r$ with $\bar{C}_{0,0} = 1$: the monopole is
carried entirely by $\mu$, which is why the model's acceleration always
contains the central term. For a center-of-mass origin the degree-1
coefficients vanish, which every shipped field satisfies. The dynamic
form factor follows from equation~\eqref{eq:gravity:norm} at
$(n,m) = (2,0)$: $N_{2,0} = \sqrt{5}$, so
$J_2 = -C_{2,0} = -\sqrt{5}\,\bar{C}_{2,0}$.

Evaluating equation~\eqref{eq:gravity:potential} through latitude and
longitude derivatives is singular at the poles, where the local
east--north basis degenerates ($1/\cos\varphi$ factors). The Pines
formulation removes the singularity by construction.

\subsection{Pines coordinates and Helmholtz polynomials}

Following Pines~\cite{pines1973}, parameterize the field point by its
radius and direction cosines,
\begin{equation}
  s = \frac{x}{r}, \qquad t = \frac{y}{r}, \qquad u = \frac{z}{r}
    = \sin\varphi,
  \label{eq:gravity:pines}
\end{equation}
and define the Helmholtz polynomials
$A_{nm}(u) = (1 - u^2)^{-m/2} P_{nm}(u)$, which are polynomials in $u$
(the $(1-u^2)^{m/2}$ factor of $P_{nm}$ is divided out). Because
$(1 - u^2)^{m/2} = \cos^m\varphi$, the longitude factors combine with
it into the complex power
\begin{equation}
  r_m(s, t) + i\, i_m(s, t) = (s + i t)^m
    = \cos^m\varphi\, e^{i m \lambda},
  \label{eq:gravity:rmim}
\end{equation}
whose real and imaginary parts obey the product recursion $r_0 = 1$,
$i_0 = 0$, $r_m = s\, r_{m-1} - t\, i_{m-1}$,
$i_m = s\, i_{m-1} + t\, r_{m-1}$ --- polynomials in $(s, t)$ with no
trigonometry, regular at the poles where $s = t = 0$. Each potential
term then factors as
\begin{equation}
  \bar{P}_{nm}(u)\left[\bar{C}_{nm}\cos m\lambda
    + \bar{S}_{nm}\sin m\lambda\right]
  = \bar{A}_{nm}(u)\, D_{nm}(s,t), \qquad
  D_{nm} = \bar{C}_{nm}\, r_m + \bar{S}_{nm}\, i_m,
  \label{eq:gravity:dnm}
\end{equation}
with $\bar{A}_{nm} = N_{nm} A_{nm}$ the normalized Helmholtz values.

\subsection{Normalized recursions}

The unnormalized Helmholtz polynomials satisfy~\cite{pines1973,
lundberg1988}
\begin{equation}
  A_{mm} = (2m-1)!!, \qquad
  A_{m+1,m} = u\,(2m+1)\,A_{mm}, \qquad
  (n-m)\,A_{nm} = u\,(2n-1)\,A_{n-1,m} - (n+m-1)\,A_{n-2,m},
  \label{eq:gravity:unnorm}
\end{equation}
and the derivative identity $\mathrm{d}A_{nm}/\mathrm{d}u =
A_{n,m+1}$. Normalized forms follow by multiplying each relation by
$N_{nm}$ and absorbing ratios of normalization factors, the approach of
Lundberg and Schutz~\cite{lundberg1988}. For the diagonal,
$A_{mm} = (2m-1)A_{m-1,m-1}$ gives
$\bar{A}_{mm} = (2m-1)\,(N_{mm}/N_{m-1,m-1})\,\bar{A}_{m-1,m-1}$; with
$N_{mm} = \sqrt{(2-\delta_{0m})(2m+1)/(2m)!}$ the ratio collapses to
\begin{equation}
  \bar{A}_{mm} =
  \sqrt{\frac{(2-\delta_{0m})}{(2-\delta_{0,m-1})}\cdot
        \frac{2m+1}{2m}}\;\bar{A}_{m-1,m-1}
  =
  \begin{cases}
    \sqrt{3}\;\bar{A}_{0,0}, & m = 1,\\[2pt]
    \sqrt{\dfrac{2m+1}{2m}}\;\bar{A}_{m-1,m-1}, & m \geq 2,
  \end{cases}
  \qquad \bar{A}_{0,0} = 1,
  \label{eq:gravity:diag}
\end{equation}
the $m = 1$ step carrying the extra $\sqrt{2}$ from the
$(2-\delta_{0m})$ factor switching on between order 0 and order 1. The
same substitution applied to the second and third of
equations~\eqref{eq:gravity:unnorm}, with
$N_{m+1,m}/N_{mm} = \sqrt{(2m+3)}/(2m+1)$,
$N_{nm}/N_{n-1,m} = \sqrt{(2n+1)(n-m)/\bigl((2n-1)(n+m)\bigr)}$, and
$N_{nm}/N_{n-2,m} =
\sqrt{(2n+1)(n-m)(n-m-1)/\bigl((2n-3)(n+m)(n+m-1)\bigr)}$, yields
\begin{equation}
  \bar{A}_{m+1,m} = u\,\sqrt{2m+3}\;\bar{A}_{mm},
  \label{eq:gravity:subdiag}
\end{equation}
\begin{equation}
  \bar{A}_{nm} =
  u\,\sqrt{\frac{(2n-1)(2n+1)}{(n-m)(n+m)}}\;\bar{A}_{n-1,m}
  \;-\;
  \sqrt{\frac{(2n+1)(n+m-1)(n-m-1)}{(2n-3)(n+m)(n-m)}}\;\bar{A}_{n-2,m},
  \label{eq:gravity:column}
\end{equation}
for $n \geq m+2$, and, from the derivative identity with
$N_{nm}/N_{n,m+1} = \sqrt{(2-\delta_{0m})(n-m)(n+m+1)/2}$,
\begin{equation}
  \frac{\mathrm{d}\bar{A}_{nm}}{\mathrm{d}u}
  = \sqrt{(n-m)(n+m+1)\cdot
      \tfrac{1}{2}(2-\delta_{0m})}\;\bar{A}_{n,m+1},
  \qquad \bar{A}_{n,m+1} \equiv 0 \text{ for } m = n.
  \label{eq:gravity:deriv}
\end{equation}

\subsection{Gradient assembly}

Write the potential as $V(r, s, t, u) = \sum_{n,m} \rho_n(r)\,
\bar{A}_{nm}(u)\, D_{nm}(s,t)$ with
$\rho_n = (\mu/r)(R/r)^n$. The gradients of the Pines coordinates are
$\nabla s = (\hat{\vv{x}} - s\hat{\vv{r}})/r$,
$\nabla t = (\hat{\vv{y}} - t\hat{\vv{r}})/r$,
$\nabla u = (\hat{\vv{z}} - u\hat{\vv{r}})/r$, and
$\nabla r = \hat{\vv{r}}$, so
\begin{equation}
  \vv{a} = \nabla V
  = a_1 \hat{\vv{x}} + a_2 \hat{\vv{y}} + a_3 \hat{\vv{z}}
  + \left[\frac{\partial V}{\partial r}
      - \bigl(s\,a_1 + t\,a_2 + u\,a_3\bigr)\right]\hat{\vv{r}},
  \qquad
  a_k \equiv \frac{1}{r}\frac{\partial V}{\partial (s,t,u)_k}.
  \label{eq:gravity:gradient}
\end{equation}
The partials of $D_{nm}$ follow from
$\partial r_m/\partial s = m\, r_{m-1}$,
$\partial i_m/\partial s = m\, i_{m-1}$,
$\partial r_m/\partial t = -m\, i_{m-1}$, and
$\partial i_m/\partial t = m\, r_{m-1}$ (differentiate
equation~\eqref{eq:gravity:rmim}), giving per-term auxiliaries
$E_{nm} = \bar{C}_{nm} r_{m-1} + \bar{S}_{nm} i_{m-1}$ and
$F_{nm} = \bar{S}_{nm} r_{m-1} - \bar{C}_{nm} i_{m-1}$. Substituting,
using $\partial \rho_n/\partial r = -(n+1)\rho_n / r$ and the identity
$s E_{nm} + t F_{nm} = D_{nm}$ (expand and regroup by
equation~\eqref{eq:gravity:rmim}), the four Pines sums are
\begin{equation}
  \begin{aligned}
    a_1 &= \sum_{n,m} \frac{\rho_n}{r}\, m\, \bar{A}_{nm}\, E_{nm},
    &\quad
    a_2 &= \sum_{n,m} \frac{\rho_n}{r}\, m\, \bar{A}_{nm}\, F_{nm},\\
    a_3 &= \sum_{n,m} \frac{\rho_n}{r}\, \bar{A}'_{nm}\, D_{nm},
    &\quad
    a_4 &= \sum_{n,m} \frac{\rho_n}{r}
      \left[(n+m+1)\,\bar{A}_{nm} + u\,\bar{A}'_{nm}\right] D_{nm},
  \end{aligned}
  \label{eq:gravity:sums}
\end{equation}
with $\bar{A}'_{nm} = \mathrm{d}\bar{A}_{nm}/\mathrm{d}u$ from
equation~\eqref{eq:gravity:deriv}, and the acceleration assembles as
\begin{equation}
  \boxed{\;\vv{a} = \bigl(a_1 - s\,a_4\bigr)\hat{\vv{x}}
        + \bigl(a_2 - t\,a_4\bigr)\hat{\vv{y}}
        + \bigl(a_3 - u\,a_4\bigr)\hat{\vv{z}}.\;}
  \label{eq:gravity:accel}
\end{equation}
Consistency check at degree 0: $\bar{A}_{0,0} = 1$, $D_{0,0} =
\bar{C}_{0,0} = 1$, $a_1 = a_2 = a_3 = 0$, $a_4 = \mu/r^2$, so
equation~\eqref{eq:gravity:accel} reduces to
$-(\mu/r^2)\,\hat{\vv{r}}$, the two-body acceleration of
Chapter~\ref{ch:twobody}. No factor anywhere divides by
$\cos\varphi$ or $\sqrt{1-u^2}$: the exact poles ($s = t = 0$,
$u = \pm 1$) are regular points, the property FR-5 selects the Pines
formulation for.

\subsection{$J_2$ secular rates (acceptance formulas)}

The first-order secular rates of the node and argument of perigee under
$J_2$, in mean elements, are~\cite{vallado2013}
\begin{equation}
  \dot{\Omega} = -\frac{3}{2}\, n\, J_2
    \left(\frac{R}{p}\right)^{\!2} \cos i,
  \qquad
  \dot{\omega} = \frac{3}{4}\, n\, J_2
    \left(\frac{R}{p}\right)^{\!2} \bigl(5\cos^2 i - 1\bigr),
  \label{eq:gravity:j2secular}
\end{equation}
with $n = \sqrt{\mu/a^3}$ and $p = a(1 - e^2)$. These are the analytic
side of the Phase~3 exit-criterion-2 gate
(\testid{GRAV-J2-SECULAR}); the numerical side is a 30-day $J_2$-only
propagation whose fitted rates must agree within \qty{0.5}{\percent}.

\section{Discretization and implementation}
\label{sec:gravity:implementation}

\begin{enumerate}
  \item \textbf{Module and traceability.} The model lives in
    \texttt{cpp/src/models/gravity.cpp}; the source comments echo the
    labels \texttt{eq:gravity:potential}, \texttt{eq:gravity:pines},
    \texttt{eq:gravity:diag}, \texttt{eq:gravity:subdiag},
    \texttt{eq:gravity:column}, \texttt{eq:gravity:deriv},
    \texttt{eq:gravity:rmim}, \texttt{eq:gravity:sums}, and
    \texttt{eq:gravity:accel} verbatim.
  \item \textbf{Precomputed recursion tables; no per-call allocation.}
    The square-root coefficients of
    equations~\eqref{eq:gravity:diag}--\eqref{eq:gravity:deriv} depend
    only on $(n, m)$; the evaluator computes them once at construction
    for the field's full degree, together with all workspace
    ($\bar{A}$ triangle, $r_m/i_m$, $\rho_n$). The per-call path is
    multiply--add only and allocation-free (D-10: no heap allocation in
    the propagation loop).
  \item \textbf{Fixed evaluation order (D-10).} Per call: the
    $\bar{A}$ triangle is filled diagonal first, then first
    sub-diagonal, then column by column in ascending $n$ for each $m$
    (the forward-stable direction~\cite{lundberg1988}); the four sums of
    equation~\eqref{eq:gravity:sums} accumulate in ascending $n$ outer,
    ascending $m$ inner order. The full triangle ($m \leq n$) is
    computed regardless of the order truncation because
    equation~\eqref{eq:gravity:deriv} reads column $m + 1$.
  \item \textbf{Tiers and truncation.} The FR-5 tiers are loop bounds:
    point-mass evaluates $n = 0$ only; $J_2$-only evaluates $n \in
    \{0, 2\}$, $m = 0$ (skipping $n = 1$ so a hypothetical non-zero
    stored degree-1 row cannot leak into the tier); full $n \times m$
    evaluates the requested band, which must lie within the stored band
    (out-of-band requests throw; Section~\ref{sec:gravity:domain}).
    Runtime truncation is arithmetically identical to load-time
    truncation --- the same terms in the same order --- which
    \testid{gravity_truncation_consistency} pins at bit level.
  \item \textbf{Pole behavior and underflow.} At $|u| \to 1$ the
    $\bar{A}_{nm}$ grow while $r_m, i_m \sim \cos^m\varphi$ decay;
    products remain bounded (Section~\ref{sec:gravity:domain}). Within
    $\sim 10^{-2}$ degrees of a pole, high-order $r_m, i_m$ underflow
    gradually to zero; the affected terms are below $10^{-300}$ of the
    total and IEEE-754 gradual underflow (guaranteed by the D-10 flags)
    keeps the result deterministic.
  \item \textbf{SRGRAV loader.} Binary read only (D-2): magic and major
    version checked first, then degree/order plausibility, tide-system
    code, and an exact file-size-versus-header check before any
    coefficient is interpreted; failures name the defect. Coefficients,
    $\mu$, and $R$ are stored and consumed as one indivisible model
    (assumption 2).
  \item \textbf{Data pipeline.} \texttt{star data fetch egm2008 |
    grgm1200a | mro120f} downloads each pinned source (SHA-256), parses
    the ICGEM gfc or PDS SHADR text in Python (D-2 boundary), truncates
    to the FR-5 degrees, and writes SRGRAV; the committed test excerpts
    under \texttt{tests/golden/gravity/} carry the same data in CSV and
    SRGRAV form with byte-anchored equality between the two
    (provenance: that directory's \texttt{manifest.toml}).
\end{enumerate}

\section{Validation evidence}
\label{sec:gravity:validation}

Table~\ref{tab:gravity:validation} lists the tests that constitute this
chapter's validation evidence. Each test identifier is machine-checked
to exist in the test tree by \texttt{scripts/lint\_chapters.py}; golden
provenance and tolerances are recorded in
\texttt{tests/golden/gravity/manifest.toml}.

\begin{table}[htbp]
  \centering
  \caption{Validation evidence for the spherical-harmonic gravity
    module.}
  \label{tab:gravity:validation}
  \begin{tabular}{lp{8.6cm}}
    \toprule
    Test ID & Acceptance basis \\
    \midrule
    \testid{GRAV-XTOOL-20} &
      Phase~3 exit criterion 1: acceleration matches independently
      synthesized pyshtools~4.14.1 values over the same committed
      excerpt coefficients at 20 states (Earth $8\times8$ and
      $20\times20$, Moon $50\times50$, Mars $20\times20$; latitudes to
      $\pm\qty{89.95}{\degree}$) to $< 10^{-12}$ relative on the
      acceleration vector (observed worst \num{2.7e-15}). \\
    \testid{GRAV-J2-SECULAR} &
      Phase~3 exit criterion 2: 30-day $J_2$-only LEO propagation
      (RKF7(8), $a = \qty{7078.137}{\kilo\metre}$, $e = 0.02$,
      $i = \qty{97.8}{\degree}$); fitted nodal and apsidal rates match
      equation~\eqref{eq:gravity:j2secular} within \qty{0.5}{\percent}
      (observed \qty{0.091}{\percent} and \qty{0.180}{\percent}). \\
    \testid{gravity_pointmass_tier} &
      Point-mass tier and point-mass factory equal the closed form
      $-\mu\hat{\vv{r}}/r^2$ ($< 10^{-15}$ relative) and each other
      (bit equality). \\
    \testid{gravity_j2_tier_closed_form} &
      $J_2$-only tier equals the closed-form central-plus-$J_2$
      acceleration at four states to $< 10^{-14}$ relative, pinning the
      $J_2 = -\sqrt{5}\,\bar{C}_{2,0}$ normalization chain. \\
    \testid{gravity_pole_regularity} &
      Evaluation exactly on both poles ($s = t = 0$, $u = \pm 1$) is
      finite, radially dominated, and continuous against a
      \qty{1}{\metre} off-axis point (bound \qty{1e-4}{\metre\per\second\squared},
      versus $\sim\qty{4e-6}{\metre\per\second\squared}$ expected from
      the gradient). \\
    \testid{gravity_truncation_consistency} &
      Runtime $8\times8$ truncation of the $20\times20$ field is
      bit-identical to a load-time-truncated field; repeated evaluation
      is bit-identical (D-10); out-of-band degree/order requests and
      $r = 0$ throw. \\
    \testid{gravity_srgrav_error_paths} &
      SRGRAV loader rejects bad magic, unsupported major version,
      truncated bodies, unknown tide codes, and missing files with
      specific errors. \\
    \testid{test_write_read_bit_exact} &
      pytest: SRGRAV write/read round trip is bit-exact and the writer
      is a pure function (identical bytes on identical input). \\
    \testid{test_csv_rebuilds_committed_srgrav_bytes} &
      pytest: each committed CSV excerpt rebuilds its committed SRGRAV
      excerpt byte-identically, so the two committed forms cannot
      drift apart. \\
    \testid{test_committed_headers_match_sources} &
      pytest: committed excerpt headers carry each source's own $GM$,
      $R$, and tide system (never \texttt{constants.hpp} values), the
      monopole is exactly 1, and degree-1 rows are zero. \\
    \bottomrule
  \end{tabular}
\end{table}

\section{References}
\label{sec:gravity:references}

This chapter relies on Pines~\cite{pines1973} for the singularity-free
formulation (direction-cosine coordinates, Helmholtz polynomials, and
the four-sum gradient assembly); on Lundberg and
Schutz~\cite{lundberg1988} for the normalized recursion approach and
its stability properties; on Vallado~\cite{vallado2013} for the
harmonic-potential development and the $J_2$ secular-rate formulas; on
Pavlis et al.~\cite{pavlis2012} for EGM2008; on Lemoine et
al.~\cite{lemoine2014} and Goossens et al.~\cite{goossens2020} for
GRGM1200A; on Konopliv et al.~\cite{konopliv2020} for MRO120F; on
Wieczorek and Meschede~\cite{wieczorek2018} for the pyshtools
cross-validation tool; and on Ince et al.~\cite{ince2019} for the ICGEM
distribution service. Full bibliographic details appear in the
bibliography.
